% =============================================================
% PREÂMBULO — adicione ao seu .tex se ainda não tiver
% =============================================================
% \usepackage{algorithm}
% \usepackage{algpseudocode}
% \usepackage{amsmath}
% \usepackage{amssymb}
% =============================================================


% =============================================================
% ALGORITMO 1 — Procedimento experimental completo
% =============================================================
\begin{algorithm}[t]
\caption{Procedimento experimental para comparação MAB vs.\ Round Robin}
\label{alg:experiment}
\begin{algorithmic}[1]
\Require Stack 5G inicializado (Open5GS Core, UERANSIM RAN, MEC Platform, gNB)
\Require Estratégias $\mathcal{S} = \{\text{MAB}, \text{RR}\}$, duração $T_{\text{run}} = 60$~min
\Statex
\State \textbf{Setup base:}
\State Inicializar 2 instâncias VidProc base $\{s_0, s_1\}$
\State Iniciar MEC Metric Catcher
\Statex
\For{cada estratégia $\sigma \in \mathcal{S}$}
    \State Limpar estado anterior (CSVs, instâncias extras)
    \State Iniciar MEC Intelligence com roteador $\sigma$
    \State Iniciar laço MAPE-K (Algoritmo~\ref{alg:mapek})
    \Statex
    \LineComment{Workload gerado em 5 fases (totaliza $T_{\text{run}}$)}
    \State \textbf{Fase 1 — Warmup} (10\,min): manter 10 UEs ativas
    \State \textbf{Fase 2 — Ramp-up} (15\,min): adicionar 1 UE a cada 10\,s até 100
    \State \textbf{Fase 3 — Peak} (15\,min): manter 100 UEs ativas
    \State \textbf{Fase 4 — Drain} (15\,min): remover 1 UE a cada 10\,s até 10
    \State \textbf{Fase 5 — Settle} (5\,min): manter 10 UEs ativas
    \Statex
    \State Arquivar métricas (rl\_input\_state.csv, rewards.csv)
    \State Encerrar Intelligence, MAPE-K e UEs
\EndFor
\Statex
\State \textbf{Análise:} computar KPIs comparativos:
\State \quad latência média, $p_{50}$, $p_{95}$, $p_{99}$
\State \quad taxa de violação de SLA $= |\{\ell : \ell > L_{\text{SLA}}\}| / N$
\State \quad \emph{instance-minutes} $= \int_0^{T_{\text{run}}} n(t)\, dt$
\State \quad número de mudanças de scaling
\end{algorithmic}
\end{algorithm}


% =============================================================
% ALGORITMO 2 — Roteamento MAB (Thompson Sampling)
% =============================================================
\begin{algorithm}[t]
\caption{Roteamento MAB com Thompson Sampling por UE}
\label{alg:mab}
\begin{algorithmic}[1]
\Require UE $u$ requisitando serviço, conjunto de braços $\mathcal{A}$ (instâncias VidProc ativas)
\Require Referência de latência $L_{\text{ref}} = 100$\,ms, dwell mínimo $T_{\text{dwell}} = 30$\,s
\Statex
\State \LineComment{Sincronizar braços com instâncias ativas, preservando estado}
\State Para braços que persistem: manter $(\alpha_a, \beta_a)$
\State Para braços novos: $(\alpha_a, \beta_a) \gets$ \textit{warm prior} (média dos atuais, cap=10)
\Statex
\State \LineComment{Atualizar estatística do braço da atribuição anterior}
\If{$u$ tem braço prévio $a_{\text{prev}}$ \textbf{e} latência medida $\ell$}
    \State $r \gets \mathbb{1}[\ell \leq L_{\text{ref}}]$ \Comment{1 se latência aceitável, 0 caso contrário}
    \State $\alpha_{a_{\text{prev}}} \gets \alpha_{a_{\text{prev}}} + r$
    \State $\beta_{a_{\text{prev}}} \gets \beta_{a_{\text{prev}}} + (1 - r)$
\EndIf
\Statex
\State \LineComment{Decidir reatribuição se dwell expirou}
\State $\Delta t \gets t_{\text{now}} - t_{\text{assign}}(u)$
\If{$\Delta t \geq T_{\text{dwell}}$ \textbf{ou} $a_{\text{prev}} \notin \mathcal{A}$}
    \For{cada braço $a \in \mathcal{A}$}
        \State $\theta_a \sim \text{Beta}(\alpha_a, \beta_a)$
    \EndFor
    \State $a^{*} \gets \arg\max_{a \in \mathcal{A}} \theta_a$
    \State Atribuir $u$ a $a^{*}$, atualizar $t_{\text{assign}}(u) \gets t_{\text{now}}$
\Else
    \State Manter $u$ em $a_{\text{prev}}$
\EndIf
\end{algorithmic}
\end{algorithm}


% =============================================================
% ALGORITMO 3 — Roteamento Round Robin (baseline)
% =============================================================
\begin{algorithm}[t]
\caption{Roteamento Round Robin (baseline)}
\label{alg:rr}
\begin{algorithmic}[1]
\Require UE $u$, conjunto de braços $\mathcal{A} = \{a_1, \ldots, a_k\}$ ordenado
\Require Dwell mínimo $T_{\text{dwell}} = 30$\,s, ponteiro circular $i$ (estado global)
\Statex
\State $\Delta t \gets t_{\text{now}} - t_{\text{assign}}(u)$
\If{$\Delta t \geq T_{\text{dwell}}$ \textbf{ou} $a_{\text{prev}} \notin \mathcal{A}$}
    \State $a^{*} \gets a_{(i \bmod k) + 1}$ \Comment{próximo na ordem cíclica}
    \State $i \gets i + 1$
    \State Atribuir $u$ a $a^{*}$, atualizar $t_{\text{assign}}(u) \gets t_{\text{now}}$
\Else
    \State Manter $u$ em $a_{\text{prev}}$
\EndIf
\end{algorithmic}
\end{algorithm}


% =============================================================
% ALGORITMO 4 — Laço MAPE-K (auto-scaling)
% =============================================================
\begin{algorithm}[t]
\caption{Laço MAPE-K para auto-scaling de instâncias VidProc}
\label{alg:mapek}
\begin{algorithmic}[1]
\Require Limites $\tau_{\text{up}} = 10\%$, $\tau_{\text{down}} = 5\%$
\Require Bounds $N_{\min} = 2$, $N_{\max} = 10$
\Require Cooldown $T_{\text{cool}} = 30$\,s, polling $T_{\text{poll}} = 5$\,s
\Statex
\State $n \gets$ \Call{InitFromRegistry}{} \Comment{maior índice ativo + 1}
\State $t_{\text{last}} \gets 0$
\Loop\ a cada $T_{\text{poll}}$ segundos
    \State \LineComment{\textbf{Monitor}: coletar métrica agregada}
    \State $c \gets \max_{i \in \mathcal{A}} \text{CPU}_i$ \Comment{CPU máxima entre instâncias}
    \Statex
    \State \LineComment{\textbf{Analyze}: decidir ação}
    \If{$t_{\text{now}} - t_{\text{last}} < T_{\text{cool}}$}
        \State $\text{ação} \gets \texttt{NO\_ACTION}$ \Comment{cooldown ativo}
    \ElsIf{$c > \tau_{\text{up}}$ \textbf{ e } $n < N_{\max}$}
        \State $\text{ação} \gets \texttt{SCALE\_UP}$
    \ElsIf{$c < \tau_{\text{down}}$ \textbf{ e } $n > N_{\min}$}
        \State $\text{ação} \gets \texttt{SCALE\_DOWN}$
    \Else
        \State $\text{ação} \gets \texttt{NO\_ACTION}$
    \EndIf
    \Statex
    \State \LineComment{\textbf{Plan + Execute}}
    \If{ação $=$ \texttt{SCALE\_UP}}
        \State Lançar instância $s_n$ via \texttt{mec\_instance\_manager}
        \If{$s_n$ aparece no \emph{service registry} em $T_{\text{verify}}$}
            \State $n \gets n + 1$; $t_{\text{last}} \gets t_{\text{now}}$
        \EndIf
    \ElsIf{ação $=$ \texttt{SCALE\_DOWN}}
        \State Encerrar instância $s_{n-1}$ via API \texttt{shut\_down\_mec\_app}
        \If{sucesso (HTTP 200)}
            \State $n \gets n - 1$; $t_{\text{last}} \gets t_{\text{now}}$
        \EndIf
    \EndIf
\EndLoop
\end{algorithmic}
\end{algorithm}


% =============================================================
% MACRO AUXILIAR (caso seu pacote algpseudocode não defina LineComment)
% =============================================================
% \algnewcommand{\LineComment}[1]{\State \(\triangleright\) #1}
%
% Coloque essa linha no preâmbulo se compilar com erro em \LineComment.
